\section{Análisis de ondas Rayleigh}
\label{sec:metodologia}

Para observar dispersión deben conservarse dos dimensiones: la evolución temporal del movimiento y la posición donde se obtuvo cada traza. Un receptor aislado describe lo ocurrido en un punto; comparar puntos separados permite reconocer cómo se desplaza la fase de una misma componente. Si $x_m$ es la posición y $t_n$ el instante de muestreo, el registro se organiza como
\begin{equation}
U[m,n]=u_z(x_m,t_n).
\label{eq:registro-xt}
\end{equation}
Cada fila es una traza temporal y el conjunto forma una matriz espacio--tiempo. Una llegada coherente aparece como una tendencia entre posiciones, siempre que se conserven la geometría y la referencia temporal.

Una primera técnica es el análisis espectral de ondas superficiales (SASW). Utiliza dos receptores y estima, para cada frecuencia, la diferencia de fase entre ambos. Su configuración es compacta, pero la fase medida admite ambigüedades de un número entero de ciclos y la escasa redundancia dificulta separar interferencias; por ello suele requerir separaciones adicionales \citep{Foti2014}.

En cambio, el análisis multicanal de ondas superficiales (MASW) extiende la observación a numerosos receptores o posiciones. La continuidad espacial permite reconocer tendencias coherentes aun en presencia de ruido y distinguir mejor contribuciones superpuestas. Desde su introducción se ha popularizado para la caracterización somera porque trabaja desde la superficie, cubre un tramo del terreno en una campaña y conserva registros que pueden reprocesarse con distintos criterios \citep{Park1998,Park1999,Foti2018}.

El costo instrumental crece con cada posición observada simultáneamente: se requiere otro transductor, acondicionamiento, conversión y una base temporal comparable. Cuando no se dispone del arreglo completo, puede sintetizarse la dimensión espacial desplazando un receptor y repitiendo la excitación bajo condiciones controladas, estrategia demostrada por \citet{LinAshlock2016}. Cada traza proviene entonces de un impacto diferente, por lo que deben registrarse tanto la posición como la excitación asociada. El diseño propuesto admite este recorrido secuencial y queda abierto a incorporar varios receptores.

La elección de MASW fija las magnitudes que debe preservar la adquisición: banda suficiente para las componentes de interés, amplitud y fase, espaciado temporal estable, referencia del impacto y geometría de cada registro. El primer elemento de esta cadena es el transductor que convierte el movimiento vertical del terreno en una señal eléctrica. La disponibilidad del geófono SM-24 y su respuesta limitada bajo la frecuencia natural originan las decisiones del sistema de adquisición descritas a continuación.
